\chapter{Forms}

Forms are where most of the frontend's interaction with the backend
happens: login, registration, and every create/edit screen for tasks. This
chapter covers the practical pattern used across all of them --- controlled
inputs backed by a single state object, validated before submit, and wired
to the API with proper loading/error handling.

\section{Controlled Inputs and a Single Form State Object}

A controlled input's value is driven entirely by React state, and every
keystroke updates that state through \texttt{onChange}. Rather than one
\texttt{useState} per field, keep the whole form as a single object and use
one generic handler keyed by the input's \texttt{name} attribute.

\begin{lstlisting}[language=JavaScript]
const [form, setForm] = useState({ email: "", password: "" });

const handleChange = (e) => {
  setForm({ ...form, [e.target.name]: e.target.value });
};

// <input name="email" value={form.email} onChange={handleChange} />
// <input name="password" value={form.password} onChange={handleChange} />
\end{lstlisting}

\begin{notebox}[NOTE]
This one handler works for any number of fields as long as each input's
\texttt{name} matches a key in the form object. Add a field to the form by
adding a key to the initial state and an input with a matching
\texttt{name} --- no new handler needed.
\end{notebox}

\section{Example 1: LoginForm}

\begin{lstlisting}[language=JavaScript]
// components/LoginForm.jsx
import { useState } from "react";
import api from "../api/axios";

function LoginForm({ onSuccess }) {
  const [form, setForm] = useState({ email: "", password: "" });
  const [errors, setErrors] = useState({});
  const [loading, setLoading] = useState(false);
  const [serverError, setServerError] = useState("");

  const handleChange = (e) => {
    setForm({ ...form, [e.target.name]: e.target.value });
  };

  const validate = () => {
    const errs = {};
    if (!form.email) errs.email = "Email is required";
    if (!form.password) errs.password = "Password is required";
    return errs;
  };

  const handleSubmit = async (e) => {
    e.preventDefault();
    const errs = validate();
    setErrors(errs);
    if (Object.keys(errs).length) return;

    setLoading(true);
    setServerError("");
    try {
      const res = await api.post("/auth/login", form);
      onSuccess(res.data.data);
      setForm({ email: "", password: "" });
    } catch (err) {
      setServerError(err.response?.data?.message || "Login failed");
    } finally {
      setLoading(false);
    }
  };

  return (
    <form onSubmit={handleSubmit} className="space-y-4">
      <div>
        <input name="email" value={form.email} onChange={handleChange} placeholder="Email" />
        {errors.email && <p className="text-red-600 text-sm">{errors.email}</p>}
      </div>
      <div>
        <input name="password" type="password" value={form.password} onChange={handleChange} placeholder="Password" />
        {errors.password && <p className="text-red-600 text-sm">{errors.password}</p>}
      </div>
      {serverError && <p className="text-red-600 text-sm">{serverError}</p>}
      <button type="submit" disabled={loading}>
        {loading ? "Logging in..." : "Log In"}
      </button>
    </form>
  );
}

export default LoginForm;
\end{lstlisting}

\section{Example 2: RegisterForm with Confirm Password}

\begin{lstlisting}[language=JavaScript]
// components/RegisterForm.jsx
import { useState } from "react";
import api from "../api/axios";

function RegisterForm() {
  const [form, setForm] = useState({
    name: "", email: "", password: "", confirmPassword: "",
  });
  const [errors, setErrors] = useState({});
  const [loading, setLoading] = useState(false);

  const handleChange = (e) => {
    setForm({ ...form, [e.target.name]: e.target.value });
  };

  const validate = () => {
    const errs = {};
    if (!form.name) errs.name = "Name is required";
    if (!/^\S+@\S+\.\S+$/.test(form.email)) errs.email = "Enter a valid email";
    if (form.password.length < 6) errs.password = "Password must be at least 6 characters";
    if (form.confirmPassword !== form.password) errs.confirmPassword = "Passwords do not match";
    return errs;
  };

  const handleSubmit = async (e) => {
    e.preventDefault();
    const errs = validate();
    setErrors(errs);
    if (Object.keys(errs).length) return;

    setLoading(true);
    try {
      const { confirmPassword, ...payload } = form;
      await api.post("/auth/register", payload);
      setForm({ name: "", email: "", password: "", confirmPassword: "" });
      setErrors({});
    } catch (err) {
      setErrors({ form: err.response?.data?.message || "Registration failed" });
    } finally {
      setLoading(false);
    }
  };

  return (
    <form onSubmit={handleSubmit} className="space-y-4">
      <input name="name" value={form.name} onChange={handleChange} placeholder="Name" />
      {errors.name && <p className="text-red-600 text-sm">{errors.name}</p>}

      <input name="email" value={form.email} onChange={handleChange} placeholder="Email" />
      {errors.email && <p className="text-red-600 text-sm">{errors.email}</p>}

      <input name="password" type="password" value={form.password} onChange={handleChange} placeholder="Password" />
      {errors.password && <p className="text-red-600 text-sm">{errors.password}</p>}

      <input name="confirmPassword" type="password" value={form.confirmPassword} onChange={handleChange} placeholder="Confirm Password" />
      {errors.confirmPassword && <p className="text-red-600 text-sm">{errors.confirmPassword}</p>}

      {errors.form && <p className="text-red-600 text-sm">{errors.form}</p>}
      <button type="submit" disabled={loading}>
        {loading ? "Creating account..." : "Register"}
      </button>
    </form>
  );
}

export default RegisterForm;
\end{lstlisting}

\begin{warnbox}[WARNING]
Never send \texttt{confirmPassword} to the backend. It exists only for
client-side validation --- strip it from the payload before the API call, as
shown with the destructuring in \texttt{handleSubmit} above.
\end{warnbox}

\section{Example 3: TaskForm --- Create/Edit Dual Mode}

The same form component handles both creating a new task and editing an
existing one. The presence of an \texttt{id} prop determines the mode: if an
\texttt{id} is passed, the form loads that task's data and submits a
\texttt{PUT}; otherwise it starts blank and submits a \texttt{POST}.

\begin{lstlisting}[language=JavaScript]
// components/TaskForm.jsx
import { useState, useEffect } from "react";
import api from "../api/axios";

function TaskForm({ id, onSaved }) {
  const isEditMode = Boolean(id);
  const [form, setForm] = useState({ title: "", description: "", status: "pending" });
  const [errors, setErrors] = useState({});
  const [loading, setLoading] = useState(false);

  useEffect(() => {
    if (!isEditMode) return;
    const loadTask = async () => {
      const res = await api.get(`/tasks/${id}`);
      const { title, description, status } = res.data.data;
      setForm({ title, description, status });
    };
    loadTask();
  }, [id, isEditMode]);

  const handleChange = (e) => {
    setForm({ ...form, [e.target.name]: e.target.value });
  };

  const validate = () => {
    const errs = {};
    if (!form.title.trim()) errs.title = "Title is required";
    return errs;
  };

  const handleSubmit = async (e) => {
    e.preventDefault();
    const errs = validate();
    setErrors(errs);
    if (Object.keys(errs).length) return;

    setLoading(true);
    try {
      if (isEditMode) {
        await api.put(`/tasks/${id}`, form);
      } else {
        await api.post("/tasks", form);
        setForm({ title: "", description: "", status: "pending" });
      }
      onSaved();
    } catch (err) {
      setErrors({ form: err.response?.data?.message || "Could not save task" });
    } finally {
      setLoading(false);
    }
  };

  return (
    <form onSubmit={handleSubmit} className="space-y-4">
      <input name="title" value={form.title} onChange={handleChange} placeholder="Title" />
      {errors.title && <p className="text-red-600 text-sm">{errors.title}</p>}

      <textarea name="description" value={form.description} onChange={handleChange} placeholder="Description" />

      <select name="status" value={form.status} onChange={handleChange}>
        <option value="pending">Pending</option>
        <option value="in-progress">In Progress</option>
        <option value="done">Done</option>
      </select>

      {errors.form && <p className="text-red-600 text-sm">{errors.form}</p>}
      <button type="submit" disabled={loading}>
        {loading ? "Saving..." : isEditMode ? "Update Task" : "Create Task"}
      </button>
    </form>
  );
}

export default TaskForm;
\end{lstlisting}

\begin{tipbox}[TIP]
This create/edit-toggle-by-prop pattern avoids writing two nearly identical
form components. The parent decides the mode simply by rendering
\texttt{<TaskForm />} for create or \texttt{<TaskForm id={task.\_id} />} for
edit.
\end{tipbox}

\section{Resetting Forms After Success}

Always reset a create form back to its initial values after a successful
submit (see \texttt{RegisterForm} and the create branch of \texttt{TaskForm}
above). Edit forms typically don't need a reset --- they either stay on the
saved values or the parent navigates away.

\whatwhyhow
{Forms use controlled inputs backed by a single state object, a generic \texttt{handleChange}, client-side validation before submit, and loading/error state around the API call.}
{Uncontrolled or scattered per-field state makes validation and reset logic inconsistent; skipping client-side validation wastes a round trip on errors that could be caught instantly.}
{Validate in a \texttt{validate()} function called from \texttt{handleSubmit} before \texttt{preventDefault}'s API call, set field errors from its return value, and reset the form object on success.}
